La progettazione dell'assistente descritto in questa tesi richiede di distinguere
il modello linguistico dal sistema che lo utilizza. Un LLM genera testo a partire
dal contesto ricevuto; un agente, invece, mantiene uno stato, sceglie strumenti,
osserva l'esito delle proprie azioni e decide come proseguire. Il capitolo
concentra la discussione sull'apprendimento statistico, sull'architettura
Transformer e sui limiti degli LLM isolati, per poi tornare alla nozione di
agente adottata nel progetto.

\section{Intelligenza artificiale e nozione di agente}

Russell e Norvig distinguono le definizioni dell'intelligenza artificiale in base
al riferimento al pensiero, al comportamento umano o alla razionalità. Per
descrivere il sistema sviluppato è sufficiente la loro prospettiva agentiva: un
agente riceve percezioni dall'ambiente e sceglie azioni \cite{russellNorvigAIMA}.
Questa definizione permette di discutere il comportamento del software senza
attribuire al modello proprietà cognitive non osservabili.

Formalmente, un agente può essere rappresentato come una funzione che trasforma
la storia delle osservazioni in un'azione. Se $p_1, p_2, \ldots, p_t$ sono le
percezioni disponibili fino al tempo $t$, l'agente calcola una politica:
\[
    \pi(p_1, p_2, \ldots, p_t) = a_t
\]
dove $a_t$ è l'azione scelta. Nell'assistente oggetto della tesi, la storia
comprende i messaggi e i risultati restituiti dalle fonti; le azioni sono il
routing verso un ramo del workflow, l'invocazione di un tool o la produzione
della risposta.

Un classificatore che riceve un'immagine e restituisce una classe rientra nei
sistemi di IA, ma non costituisce necessariamente un agente: non interagisce con
l'ambiente nel tempo e non conserva un obiettivo persistente. Nell'agente queste
due proprietà sono centrali, perché ogni azione modifica le informazioni sulle
quali verrà presa la decisione successiva.

\section{Apprendimento statistico e reti profonde}

\subsection{Dai sistemi simbolici all'apprendimento dai dati}

La ricerca degli anni Cinquanta adottò soprattutto rappresentazioni simboliche:
fatti, predicati e regole di inferenza venivano specificati esplicitamente. Da
questa impostazione derivano, tra l'altro, la pianificazione classica, la
dimostrazione automatica di teoremi e i sistemi esperti. Il metodo resta efficace
nei domini formalizzabili, ma richiede un aggiornamento manuale quando cambiano
regole e conoscenza del dominio.

Il machine learning affronta diversamente il problema: il programmatore definisce
una famiglia di modelli e una procedura di apprendimento, mentre i parametri sono
stimati a partire dagli esempi. Sistemi simbolici e modelli appresi non
corrispondono comunque a fasi mutuamente esclusive; nell'architettura descritta
più avanti, le decisioni probabilistiche dell'LLM convivono con regole e schemi
di validazione scritti nel codice.

\subsection{Machine learning}

Nel machine learning una funzione sconosciuta viene approssimata a partire dai
dati osservati. Dato un insieme di esempi:
\[
    D = \{(x_i, y_i)\}_{i=1}^{n}
\]
si cerca una funzione $f_\theta$ parametrizzata da $\theta$ tale che $f_\theta(x_i)$ sia vicino a $y_i$. In classificazione, $x_i$ può essere un'immagine e $y_i$ una classe. In regressione, $y_i$ può essere un valore numerico. In apprendimento non supervisionato, l'obiettivo non è predire etichette, ma scoprire struttura nei dati: cluster, rappresentazioni compatte, anomalie o fattori latenti.

Rispetto a un programma tradizionale, parte del comportamento non è specificata
direttamente nel codice. Il programmatore definisce architettura, procedura di
addestramento e funzione obiettivo; il risultato dipende anche dai dati e dai
parametri appresi.

Un modello supervisionato viene allenato minimizzando una loss:
\[
    \theta^* = \arg\min_\theta \frac{1}{n} \sum_{i=1}^{n} L(f_\theta(x_i), y_i)
\]
dove $L$ misura l'errore tra predizione e valore atteso. L'architettura di
$f_\theta$, i dati e la funzione di perdita dipendono dal problema, ma i
parametri vengono in ogni caso aggiornati in funzione degli esempi osservati.

Il machine learning ha reso possibile trattare problemi in cui scrivere regole esplicite era troppo costoso o fragile. Tuttavia, molti modelli classici richiedevano ancora una fase di feature engineering: un esperto doveva decidere quali caratteristiche estrarre dai dati grezzi. In visione artificiale, per esempio, prima delle reti profonde si progettavano descrittori di bordi, texture o forme. In elaborazione del linguaggio naturale si usavano n-grammi, dizionari, feature sintattiche e rappresentazioni sparse.

\subsection{Deep learning}

Il deep learning estende il machine learning usando reti neurali con molti
strati, capaci di apprendere rappresentazioni gerarchiche. Ogni strato trasforma
l'input in una rappresentazione intermedia sulla quale opera lo strato
successivo. In una rete per immagini, per esempio, i primi livelli possono
rilevare pattern locali e quelli finali concetti più vicini alle classi di
interesse.

LeCun, Bengio e Hinton descrivono il deep learning come una famiglia di metodi in
cui modelli composti da più livelli apprendono rappresentazioni con diversi
livelli di astrazione \cite{lecunDeepLearning}. Le trasformazioni utili al compito
vengono apprese dai dati, riducendo il ricorso a feature progettate manualmente.

La disponibilità di dataset estesi e GPU, insieme ai progressi nelle architetture
e nelle procedure di ottimizzazione, ha reso praticabile l'addestramento di reti
sempre più profonde.

Nel 2012 AlexNet ottenne su ImageNet un errore di classificazione sensibilmente
inferiore rispetto agli altri sistemi presentati nella stessa competizione
\cite{krizhevskyImageNet}. Il risultato contribuì alla diffusione delle reti
profonde nella visione artificiale e mostrò l'effetto congiunto di dataset estesi,
GPU e architetture adatte al compito.

Una parte della complessità si sposta così dalle condizioni scritte nel codice ai
parametri: il comportamento dipende congiuntamente da dati, architettura e
procedura di addestramento.

\section{Modelli linguistici e architettura Transformer}

\subsection{Modellazione del linguaggio}

Il linguaggio naturale presenta ambiguità, dipendenze dal contesto e riferimenti a
conoscenze non espresse nella frase. I sistemi di NLP hanno a lungo combinato
regole linguistiche, feature manuali e modelli statistici; il deep learning ha
poi favorito rappresentazioni dense e reti neurali sequenziali.

Un modello linguistico assegna probabilità a sequenze di token. Se una frase viene rappresentata come $x_1, x_2, \ldots, x_n$, un modello autoregressivo stima:
\[
    P(x_1, \ldots, x_n) = \prod_{t=1}^{n} P(x_t \mid x_1, \ldots, x_{t-1})
\]
L'addestramento consiste quindi nella predizione del token successivo dato il
contesto precedente. Pur essendo formulato come un compito locale, questo
obiettivo richiede di rappresentare regolarità grammaticali e semantiche, oltre
alle relazioni fattuali presenti nei dati di addestramento.

Il tokenizer converte il testo in unità discrete, che possono coincidere con una
parola, una sua parte o un segno di punteggiatura. A ciascuna unità è associato
un identificativo numerico, trasformato poi in un embedding sul quale opera il
modello.

\subsection{Transformer e attenzione}

L'architettura di riferimento per gli LLM è il Transformer, introdotto da Vaswani
et al. in \emph{Attention Is All You Need} \cite{vaswaniAttention}. I precedenti
modelli sequenziali usavano spesso reti ricorrenti o convoluzionali. Le reti
ricorrenti elaborano la sequenza un elemento alla volta e mantengono uno stato
nascosto, con minori possibilità di parallelizzare l'addestramento e difficoltà
nel rappresentare dipendenze molto lunghe.

Il Transformer usa invece il meccanismo di self-attention. Ogni token costruisce una rappresentazione guardando gli altri token del contesto e pesandoli in base alla loro rilevanza. In forma semplificata, l'attenzione usa tre matrici: query $Q$, key $K$ e value $V$. Il punteggio tra un token e un altro viene calcolato confrontando query e key; i value vengono poi combinati con pesi normalizzati:
\[
    \mathrm{Attention}(Q,K,V) = \mathrm{softmax}\left(\frac{QK^T}{\sqrt{d_k}}\right)V
\]
Il calcolo mette in relazione direttamente posizioni diverse della sequenza,
senza comprimere tutto il prefisso in un singolo stato ricorrente. Le diverse
teste di attenzione apprendono relazioni differenti in parallelo. Questa
struttura facilita la parallelizzazione dell'addestramento e si è dimostrata
adatta a scalare con dati e modelli di grandi dimensioni.

\subsection{Modelli linguistici di grandi dimensioni}

Con LLM si indica un modello linguistico neurale addestrato su collezioni testuali
estese, in genere basato su Transformer. Il termine ``grande'' può riferirsi al
numero di parametri, ai dati e al calcolo impiegati o alla lunghezza del contesto.

I parametri sono i valori numerici appresi, organizzati principalmente in matrici
di pesi e vettori di bias. GPT-3, per esempio, è stato presentato come un modello
autoregressivo da 175 miliardi di parametri \cite{brownGPT3}. La sola dimensione
non misura la qualità di un modello: architettura, dati e addestramento incidono
sul risultato, mentre l'aumento dei parametri comporta maggiori costi di memoria
e calcolo.

Un parametro non corrisponde a una frase memorizzata, ma è un coefficiente che
contribuisce a molte predizioni. L'informazione appresa è distribuita tra i
pesi, mentre un database conserva record espliciti e li recupera senza
generarli. La distinzione è rilevante per questo lavoro: i dati su eventi e
acquisizioni satellitari devono provenire dalle fonti interrogate dal sistema,
non dalla memoria parametrica del modello.

Durante l'inferenza, il modello riceve un prompt, cioè una sequenza di input, e produce una distribuzione di probabilità sul token successivo. La generazione procede iterativamente: il token scelto viene aggiunto al contesto e il processo continua. Parametri di decoding come temperatura, top-$p$ o penalità di ripetizione modificano il modo in cui si campiona dalla distribuzione, influenzando determinismo e varietà dell'output.

\subsection{Pre-addestramento, adattamento e limiti degli LLM isolati}

Il ciclo tipico di un LLM moderno comprende il pre-addestramento su grandi
collezioni di testo e, spesso, un successivo adattamento alle istruzioni. Nel
contesto applicativo di questa tesi, l'uso del solo modello lascia irrisolti tre
problemi:
\begin{itemize}
    \item \textbf{Conoscenza statica}: i parametri riflettono il training set e non si aggiornano automaticamente;
    \item \textbf{Allucinazione}: il modello può generare testo plausibile ma falso;
    \item \textbf{Assenza di effetti esterni}: un modello che produce solo testo non può eseguire comandi, leggere documenti aggiornati o interagire con il mondo.
\end{itemize}

Fonti aggiornate, memoria e azioni devono quindi essere fornite da componenti
esterni al modello.

\section{Agenti basati su LLM}

Nelle architetture considerate in questa tesi, un agente IA è un sistema
software che usa un LLM per scegliere e coordinare azioni. Attorno al modello è
implementato un ciclo di controllo che riceve un obiettivo, osserva il contesto,
invoca strumenti esterni e aggiorna lo stato in base ai risultati ottenuti.

La letteratura recente ha studiato esplicitamente l'alternanza tra ragionamento
e azione. ReAct combina tracce di ragionamento e invocazioni di strumenti nello
stesso processo decisionale \cite{yaoReAct}; Toolformer affronta invece il
problema di apprendere quando e come richiamare API esterne
\cite{schickToolformer}. Il progetto qui presentato non addestra un nuovo modello: applica questi
principi a un workflow LangGraph controllato (\emph{Copernicus Data Workspace}),
nel quale le decisioni dell'LLM sono validate da schemi e regole prima
dell'esecuzione sulle sorgenti.

Il passaggio dai sistemi simbolici agli agenti contemporanei non va inteso come
una sostituzione lineare: regole esplicite e modelli appresi convivono nelle
stesse architetture. In questa prospettiva, il modello linguistico è una
componente dell'agente e non la sua intera implementazione. Nel prototipo
azioni, memoria Redis, artefatti dashboard e controlli restano affidati a
componenti software separati. Il capitolo seguente esamina MCP, il protocollo
adottato per collegare alcuni di questi componenti a GDACS e CEMS.
